\subsection{ALMA and Hubble Astrometry} \label{sec:astrometry}

    The position of the LBG was originally determined by the HFF \HST{} images which are aligned to the existing CLASH catalogs \citep{Postman12}; the CLASH astrometry was based on Subaru's Suprime-Cam catalogs which are registered onto the Two-Micron All Sky Survey (2MASS) frame.  We find, however, that the optical-to-NIR astrometry does not fully match the International Celestial Reference System (ICRS), on which ALMA relies.  In order to correct the astrometry of the \HST{} images, we use four positions of objects (3 stars and 1 cluster elliptical) accurately measured in the \Gaia{} first data release (DR1) catalog \citep{Gaia16a, Gaia16b}. We also compare the positions of the three quasars used as ALMA phase calibrators, J0348$-$2749, J0453$-$2807 and J0522$-$3627, which are determined by the \Gaia{} DR1 catalog and by phase solutions from our ALMA calibration.
    We find that the relative offsets between the ALMA and \Gaia{} coordinates are typically $< 10$ mas. All of the \HST{} images are corrected for astrometry on the basis of the \Gaia{} coordinates using the \textsc{iraf} \citep{Tody93} task, \texttt{ccsetwcs}, confirming those two frames coincide with each other down to the accuracy of $\lesssim 30$ mas. The resulting centroid of the LBG in the ICRS coordinate is at $\rm (\alpha_{ICRS}, \delta_{ICRS}) = (04^h 16^m 09\fs 415, -24\arcdeg 05' 35\farcs54)$. We hereafter use this coordinate as the formal position of \target{}.


